\section*{अध्याय \ref{ch:g7:atoms-and-molecules} --- परमाणु और अणु}

\begin{solution}{exo:g7:atoms-and-molecules:1}
\omterm{def:g7:atoms-and-molecules:atom}{परमाणु} उन बहुत ही छोटे कणों में से एक है जिनसे सारा द्रव्य बना है
(इनके लगभग सौ प्रकार हैं)। \omterm{def:g7:atoms-and-molecules:molecule}{अणु} \omterm{def:g7:atoms-and-molecules:atom}{परमाणुओं} का मज़बूती से जुड़ा एक समूह है,
किसी पदार्थ का वह सबसे छोटा टुकड़ा जो अब भी वही
पदार्थ है।
\end{solution}

\begin{solution}{exo:g7:atoms-and-molecules:2}
\ce{H}, \ce{C}, \ce{N}, \ce{O}, \ce{Na}, \ce{Fe}।
\end{solution}

\begin{solution}{exo:g7:atoms-and-molecules:3}
1 कार्बन \omterm{def:g7:atoms-and-molecules:atom}{परमाणु} और 2 ऑक्सीजन \omterm{def:g7:atoms-and-molecules:atom}{परमाणु}: कुल 3 \omterm{def:g7:atoms-and-molecules:atom}{परमाणु}।
\end{solution}

\begin{solution}{exo:g7:atoms-and-molecules:4}
पानी, ऑक्सीजन, मेथेन, नाइट्रोजन।
\end{solution}

\begin{solution}{exo:g7:atoms-and-molecules:5}
\ce{NO2}।
\end{solution}

\begin{solution}{exo:g7:atoms-and-molecules:6}
\ce{O2} एक \omterm{def:g7:atoms-and-molecules:molecule}{अणु} है जो दो जुड़े हुए ऑक्सीजन \omterm{def:g7:atoms-and-molecules:atom}{परमाणुओं} से बना है; \ce{2O} दो
अलग-अलग ऑक्सीजन \omterm{def:g7:atoms-and-molecules:atom}{परमाणु} हैं। \ce{CO} एक कार्बन और एक ऑक्सीजन \omterm{def:g7:atoms-and-molecules:atom}{परमाणु} का
\omterm{def:g7:atoms-and-molecules:molecule}{अणु} है (दो बड़े अक्षर, दो \omterm{def:g7:atoms-and-molecules:symbol}{प्रतीक}); \ce{Co} कोबाल्ट का एक \omterm{def:g7:atoms-and-molecules:atom}{परमाणु} है
(एक \omterm{def:g7:atoms-and-molecules:symbol}{प्रतीक}, पहले बड़ा फिर छोटा अक्षर)।
\end{solution}

\begin{solution}{exo:g7:atoms-and-molecules:7}
\ce{3H2O}: $3\times 2 = 6$ हाइड्रोजन \omterm{def:g7:atoms-and-molecules:atom}{परमाणु} और $3\times 1 = 3$ ऑक्सीजन
\omterm{def:g7:atoms-and-molecules:atom}{परमाणु}। \ce{5O2}: कोई हाइड्रोजन नहीं, $5\times 2 = 10$ ऑक्सीजन \omterm{def:g7:atoms-and-molecules:atom}{परमाणु}।
\end{solution}

\begin{solution}{exo:g7:atoms-and-molecules:8}
काला कार्बन है, लाल ऑक्सीजन: एक कार्बन, दो ऑक्सीजन, \ce{CO2},
कार्बन डाइऑक्साइड।
\end{solution}

\begin{solution}{exo:g7:atoms-and-molecules:9}
डाल्टन का पानी \ce{HO} था। असली \omterm{def:g7:atoms-and-molecules:molecule}{अणु}, \ce{H2O}, में एक हाइड्रोजन
\omterm{def:g7:atoms-and-molecules:atom}{परमाणु} ज़्यादा है: डाल्टन के \omterm{def:g7:atoms-and-molecules:molecule}{अणु} में एक \omterm{def:g7:atoms-and-molecules:atom}{परमाणु} कम था।
\end{solution}

\begin{solution}{exo:g7:atoms-and-molecules:10}
शुद्ध पानी \omterm{def:g6:pure-substances-and-mixtures:mixture}{मिश्रण} नहीं है: उसके सभी \omterm{def:g7:atoms-and-molecules:molecule}{अणु} एक जैसे \ce{H2O}
\omterm{def:g7:atoms-and-molecules:molecule}{अणु} हैं। \omterm{def:g4:air-a-mixture-of-gases:air}{वायु} एक \omterm{def:g6:pure-substances-and-mixtures:mixture}{मिश्रण} है: उसमें नाइट्रोजन के \omterm{def:g7:atoms-and-molecules:molecule}{अणु}, ऑक्सीजन के
\omterm{def:g7:atoms-and-molecules:molecule}{अणु}, आर्गन \omterm{def:g7:atoms-and-molecules:atom}{परमाणु} और कार्बन डाइऑक्साइड के कुछ \omterm{def:g7:atoms-and-molecules:molecule}{अणु} हैं।
\end{solution}

\begin{solution}{exo:g7:atoms-and-molecules:11}
एक करोड़ \omterm{def:g7:atoms-and-molecules:atom}{परमाणु} \qty{1}{mm} बनाते हैं, इसलिए \qty{1}{cm} $= \qty{10}{mm}$
के लिए $10 \times$ एक करोड़ $=$ दस करोड़ \omterm{def:g7:atoms-and-molecules:atom}{परमाणु} चाहिए।
$\qty{100}{m} = \num{100000}$~mm के लिए $\num{100000}\times$ एक करोड़
$=$ एक लाख करोड़ \omterm{def:g7:atoms-and-molecules:atom}{परमाणु} चाहिए।
\end{solution}

\begin{solution}{exo:g7:atoms-and-molecules:12}
$6 + 12 + 6 = 24$ \omterm{def:g7:atoms-and-molecules:atom}{परमाणु}। हाइड्रोजन: $\frac{12}{24} = \qty{50}{\percent}$।
कार्बन: $\frac{6}{24} = \qty{25}{\percent}$ (और बाकी
\qty{25}{\percent} ऑक्सीजन)।
\end{solution}

\begin{solution}{pb:g7:atoms-and-molecules:1}
\textbf{1.} नाइट्रोजन \ce{N2}, 2 \omterm{def:g7:atoms-and-molecules:atom}{परमाणु}; ऑक्सीजन \ce{O2}, 2 \omterm{def:g7:atoms-and-molecules:atom}{परमाणु}; कार्बन
डाइऑक्साइड \ce{CO2}, 3 \omterm{def:g7:atoms-and-molecules:atom}{परमाणु}।

\textbf{2.} आर्गन \omterm{def:g7:atoms-and-molecules:atom}{परमाणु} जुड़कर \omterm{def:g7:atoms-and-molecules:molecule}{अणु} नहीं बनाते: आर्गन अकेले
\omterm{def:g7:atoms-and-molecules:atom}{परमाणुओं} से बनी है।

\textbf{3.} नाइट्रोजन: $78 \times 2 = 156$ \omterm{def:g7:atoms-and-molecules:atom}{परमाणु}। ऑक्सीजन: $21\times 2 =
42$ \omterm{def:g7:atoms-and-molecules:atom}{परमाणु}।

\textbf{4.} $156 + 42 + 1 = 199$ \omterm{def:g7:atoms-and-molecules:atom}{परमाणु}।

\textbf{5.} $\frac{156}{199} \approx 0.78$, लगभग \qty{78}{\percent}।

\textbf{6.} ऑक्सीजन के \omterm{def:g7:atoms-and-molecules:molecule}{अणुओं} में: $16\times 2 = 32$; कार्बन डाइऑक्साइड के
\omterm{def:g7:atoms-and-molecules:molecule}{अणुओं} में: $4 \times 2 = 8$; कुल $32 + 8 = 40$।

\textbf{7.} $42 - 40 = 2$ ऑक्सीजन \omterm{def:g7:atoms-and-molecules:atom}{परमाणु} कम हैं: ऑक्सीजन के एक \omterm{def:g7:atoms-and-molecules:molecule}{अणु}
जितने। वे पानी के \omterm{def:g7:atoms-and-molecules:molecule}{अणुओं}, \ce{H2O}, में गए हैं, जिन्हें शरीर अपने भोजन की
हाइड्रोजन से बनाता है; वह पानी आंशिक रूप से साँस में जलवाष्प के रूप में
निकलता है (जिसे इन आँकड़ों में हटा दिया गया है)।

\textbf{8.} अंदर ली गई: 0। बाहर छोड़ी गई: 4 (हर \ce{CO2} में एक)।

\textbf{9.} शरीर द्वारा इस्तेमाल किए गए भोजन से: शर्कराओं और वसाओं में कार्बन
\omterm{def:g7:atoms-and-molecules:atom}{परमाणु} होते हैं, और शरीर उन्हें बाहर छोड़ी गई कार्बन डाइऑक्साइड में मुक्त करता है।

\textbf{10.} $21 - 16 = 5$ ऑक्सीजन \omterm{def:g7:atoms-and-molecules:molecule}{अणु} ग़ायब हुए हैं और कार्बन
डाइऑक्साइड के 4 \omterm{def:g7:atoms-and-molecules:molecule}{अणु} प्रकट हुए हैं।

\textbf{11.} नाइट्रोजन: तीन छड़ों से जुड़ी दो नीली गेंदें। ऑक्सीजन: दो
छड़ों से जुड़ी दो लाल गेंदें। कार्बन डाइऑक्साइड: दो लाल गेंदों के बीच एक काली
गेंद, हर लाल गेंद उससे दो छड़ों से जुड़ी, एक सीधी रेखा में। आर्गन: एक
अकेली हल्की नीली गेंद।

\textbf{12.} $156 + 32 + 12 + 1 = 201$ \omterm{def:g7:atoms-and-molecules:atom}{परमाणु} (कार्बन डाइऑक्साइड के 4
\omterm{def:g7:atoms-and-molecules:molecule}{अणुओं} में 12 \omterm{def:g7:atoms-and-molecules:atom}{परमाणु} हैं)। यह अंदर ली गई \omterm{def:g4:air-a-mixture-of-gases:air}{वायु} से $201 - 199 = 2$ \omterm{def:g7:atoms-and-molecules:atom}{परमाणु}
ज़्यादा है: भोजन से मिले 4 कार्बन \omterm{def:g7:atoms-and-molecules:atom}{परमाणु}, घटाओ पानी के साथ निकले 2 ऑक्सीजन
\omterm{def:g7:atoms-and-molecules:atom}{परमाणु}।
\end{solution}
